% Generated draft: requires mathematical and source review.
\documentclass[11pt]{article}
\usepackage[margin=25mm]{geometry}
\usepackage{fontspec}
\setmainfont{Latin Modern Roman}
\usepackage{amsmath,amssymb,mathtools,unicode-math}
\setmathfont{Latin Modern Math}
\usepackage{enumitem,needspace,xurl,hyperref}
\setlength{\parindent}{0pt}
\setlength{\parskip}{5pt}
\setlength{\emergencystretch}{3em}
\newcommand{\N}{\mathbb{N}}
\newcommand{\Z}{\mathbb{Z}}
\newcommand{\Q}{\mathbb{Q}}
\newcommand{\R}{\mathbb{R}}
\newcommand{\C}{\mathbb{C}}
\newcommand{\F}{\mathbb{F}}
\newcommand{\A}{\mathbb{A}}
\newcommand{\PP}{\mathbb{P}}
\newcommand{\E}{\mathbb{E}}
\newcommand{\eps}{\varepsilon}
\newcommand{\dd}{\,\mathrm{d}}
\begin{document}
\section*{UnsolvedMath ID: 800722}
\textbf{Source problem code:} catalogue entry
\subsection{Definitions and mathematical statement}
Let \(F\) be an arbitrary field and \(G\) a finite group. A central division algebra over \(F\) is a finite-dimensional associative division algebra \(A\) whose center is \(F\); its degree is \(\deg(A)=\sqrt{\dim_F A}\). A \(G\)-crossed product division algebra over \(F\) means such an algebra \(A\) containing a maximal commutative subfield \(L\) for which \(L/F\) is Galois and \(\operatorname{Gal}(L/F)\cong G\). In particular, \([L:F]=\deg(A)=|G|\).

The group \(G\) is called \emph{\(F\)-admissible} if such an algebra exists. Determine, for an arbitrary field \(F\), precisely which finite groups \(G\) are \(F\)-admissible; equivalently, characterize the pairs \((F,G)\) for which there exist a finite-dimensional central division \(F\)-algebra \(A\) and a subfield \(L\subseteq A\) satisfying
\[
L/F\text{ is Galois},\qquad \operatorname{Gal}(L/F)\cong G,\qquad [L:F]=\deg(A)=|G|.
\]
The supplied formulation imposes no restriction on the characteristic or other properties of \(F\).
\subsection{Short English statement}
For each field \(F\), classify the finite groups that occur as the Galois group of a maximal subfield of a finite-dimensional central division algebra over \(F\). The field extension and division algebra are both allowed to vary.
\subsection{Sources}
\begin{itemize}
\item Local UnsolvedMath registry, permanent ID \(800722\), source code \texttt{catalogue entry}, titled \emph{Open Problems on Central Simple Algebras --- Problem 37}. The supplied record attributes the problem to Auel, Brussel, Garibaldi, and Vishne, with proposed year 2010.
\item Auel, Brussel, Garibaldi, and Vishne, \emph{Open Problems on Central Simple Algebras} (2010), Problem 37. Supplied source link: \url{https://arxiv.org/abs/1006.3304}. The record reports extraction from arXiv LaTeX and access on 2026-10-16; these provenance details are supplied and have not been independently verified.
\item The local record labels the problem \emph{open}. This is a supplied registry status, not an independently verified assessment of the literature.
\end{itemize}
\end{document}
